\section{1. Información general del
proyecto}\label{informaciuxf3n-general-del-proyecto}

\subsection{1.1 Integrantes y roles}\label{integrantes-y-roles}

{\def\LTcaptype{none} % do not increment counter
\begin{longtable}[]{@{}
  >{\raggedright\arraybackslash}p{(\linewidth - 8\tabcolsep) * \real{0.2456}}
  >{\raggedright\arraybackslash}p{(\linewidth - 8\tabcolsep) * \real{0.1282}}
  >{\raggedright\arraybackslash}p{(\linewidth - 8\tabcolsep) * \real{0.2012}}
  >{\raggedright\arraybackslash}p{(\linewidth - 8\tabcolsep) * \real{0.1686}}
  >{\raggedright\arraybackslash}p{(\linewidth - 8\tabcolsep) * \real{0.2564}}@{}}
\toprule\noalign{}
\begin{minipage}[b]{\linewidth}\raggedright
\textbf{Nombre completo}
\end{minipage} & \begin{minipage}[b]{\linewidth}\raggedright
\textbf{Matrícula}
\end{minipage} & \begin{minipage}[b]{\linewidth}\raggedright
\textbf{Usuario GitHub}
\end{minipage} & \begin{minipage}[b]{\linewidth}\raggedright
\textbf{Rol(es) y sprints}
\end{minipage} & \begin{minipage}[b]{\linewidth}\raggedright
\textbf{Áreas principales de contribución}
\end{minipage} \\
\midrule\noalign{}
\endhead
\bottomrule\noalign{}
\endlastfoot
Ángel Yael Avendaño Guzmán & 0224170023 & AyagDan & Product Owner
(S0--S6) & Enlace con el cliente, levantamiento de requerimientos en
entrevistas, priorización del Product Backlog y aceptación de
entregables. \\
Abdías Daniel Chávez Cruz & 0224170024 & absdan & Scrum Master /
Tech Lead (S0--S6) & Facilitación de Scrum, arquitectura del backend en
FastAPI, diseño de la base de datos relacional y gestión técnica de Pull
Requests. \\
Dulce María Estudillo Rosado & 0224170028 & DulceRosado & QA Lead /
Developer (S0--S6) & Diseño del plan de pruebas, ejecución de casos de
aceptación, validación de calidad y realización del QA cruzado al otro
equipo. \\
Isis Alhelí González Maldonado & 0224170030 & \ruta{IsisAlheliGonzalezMaldonado} &
Developer (S0--S6) & Redacción de la especificación de requisitos de
software (SRS), historias de usuario y modelado de procesos del
negocio. \\
Isis Yamilet Meléndez Hernández & 0224170037 & \ruta{IsisMelendezHernandez} &
Developer (S0--S6): especialista en interfaz, datos y documentación &
Integración y compilación del Documento de Ingeniería (migración a
LaTeX/Word), configuración de flujos de CI en GitHub Actions, despliegue
y apoyo en el desarrollo de los componentes de la interfaz (PWA). \\
\end{longtable}
}

\subsection{1.2 Enlaces del proyecto}\label{enlaces-del-proyecto}

{\def\LTcaptype{none} % do not increment counter
\begin{longtable}[]{@{}
  >{\raggedright\arraybackslash}p{(\linewidth - 2\tabcolsep) * \real{0.3631}}
  >{\raggedright\arraybackslash}p{(\linewidth - 2\tabcolsep) * \real{0.6369}}@{}}
\toprule\noalign{}
\begin{minipage}[b]{\linewidth}\raggedright
\textbf{Recurso}
\end{minipage} & \begin{minipage}[b]{\linewidth}\raggedright
\textbf{Enlace}
\end{minipage} \\
\midrule\noalign{}
\endhead
\bottomrule\noalign{}
\endlastfoot
Repositorio & \url{https://github.com/AyagDan/SistemaDeServiciosYCitas} \\
Tablero (GitHub Projects / Jira) & Historias como issues del repositorio:
\url{https://github.com/AyagDan/SistemaDeServiciosYCitas/issues} \\
Prototipo de UI (Figma / Penpot) & Se agrega con los prototipos de baja
fidelidad (4.10, Review 2). \\
Ambiente de pruebas (staging) & Desde el Sprint 4. \\
Release v0.1.0 & Se publica en GitHub Releases en la Demo final
(Sprint 6). \\
Canal de comunicación del equipo & Grupo de WhatsApp del equipo
(mensajes y llamadas); reuniones presenciales en la escuela. \\
\end{longtable}
}

\subsection{1.3 Resumen ejecutivo}\label{resumen-ejecutivo}

Servicio Automotriz Hermanos Ruiz opera hoy con cuadernos, órdenes en
papel y llamadas. Esto provoca disputas por cobros que el cliente dice
no haber autorizado, reclamos por daños sin evidencia del estado en que
llegó el vehículo, muchas llamadas para preguntar por el avance y poca
visibilidad de qué vehículo ocupa cada una de las 4 bahías.

El Sistema de Órdenes de Servicio y Citas digitalizará el ciclo de
atención: la cita, la recepción del vehículo con checklist y
fotografías, el diagnóstico y la asignación de bahía, la cotización que
el cliente autoriza desde su celular con un código de verificación, y el
cobro y la entrega. El cliente consultará el estado de su vehículo con un
enlace seguro, sin instalar aplicaciones.

Estado actual (v\versiondoc): además de la visión y la viabilidad
(sección 3), ya está en borrador la especificación de requisitos
(sección 4): el glosario con los términos del taller, los usuarios, los
requisitos funcionales, las 30 historias de usuario con su prioridad, los
requisitos de calidad medibles y las reglas de negocio, incluidos los
procedimientos para cuando falle la luz o internet. En esta Review se
recorre con el cliente para validarla; la versión firmada se entrega en
la Review 2. La Fase 2 agregará reportes gerenciales, consulta de
existencias y de garantías, y el tablero simplificado de bahías.
